% Options for packages loaded elsewhere
\PassOptionsToPackage{unicode}{hyperref}
\PassOptionsToPackage{hyphens}{url}
%
\documentclass[
]{article}
\usepackage{amsmath,amssymb}
\usepackage{iftex}
\ifPDFTeX
  \usepackage[T1]{fontenc}
  \usepackage[utf8]{inputenc}
  \usepackage{textcomp} % provide euro and other symbols
\else % if luatex or xetex
  \usepackage{unicode-math} % this also loads fontspec
  \defaultfontfeatures{Scale=MatchLowercase}
  \defaultfontfeatures[\rmfamily]{Ligatures=TeX,Scale=1}
\fi
\usepackage{lmodern}
\ifPDFTeX\else
  % xetex/luatex font selection
\fi
% Use upquote if available, for straight quotes in verbatim environments
\IfFileExists{upquote.sty}{\usepackage{upquote}}{}
\IfFileExists{microtype.sty}{% use microtype if available
  \usepackage[]{microtype}
  \UseMicrotypeSet[protrusion]{basicmath} % disable protrusion for tt fonts
}{}
\makeatletter
\@ifundefined{KOMAClassName}{% if non-KOMA class
  \IfFileExists{parskip.sty}{%
    \usepackage{parskip}
  }{% else
    \setlength{\parindent}{0pt}
    \setlength{\parskip}{6pt plus 2pt minus 1pt}}
}{% if KOMA class
  \KOMAoptions{parskip=half}}
\makeatother
\usepackage{xcolor}
\usepackage[margin=1in]{geometry}
\usepackage{graphicx}
\makeatletter
\def\maxwidth{\ifdim\Gin@nat@width>\linewidth\linewidth\else\Gin@nat@width\fi}
\def\maxheight{\ifdim\Gin@nat@height>\textheight\textheight\else\Gin@nat@height\fi}
\makeatother
% Scale images if necessary, so that they will not overflow the page
% margins by default, and it is still possible to overwrite the defaults
% using explicit options in \includegraphics[width, height, ...]{}
\setkeys{Gin}{width=\maxwidth,height=\maxheight,keepaspectratio}
% Set default figure placement to htbp
\makeatletter
\def\fps@figure{htbp}
\makeatother
\setlength{\emergencystretch}{3em} % prevent overfull lines
\providecommand{\tightlist}{%
  \setlength{\itemsep}{0pt}\setlength{\parskip}{0pt}}
\setcounter{secnumdepth}{-\maxdimen} % remove section numbering
\usepackage[spanish]{babel}
\usepackage{amsmath}
\usepackage{tikz}
\ifLuaTeX
  \usepackage{selnolig}  % disable illegal ligatures
\fi
\usepackage{bookmark}
\IfFileExists{xurl.sty}{\usepackage{xurl}}{} % add URL line breaks if available
\urlstyle{same}
\hypersetup{
  hidelinks,
  pdfcreator={LaTeX via pandoc}}

\author{}
\date{\vspace{-2.5em}}

\begin{document}

\section{Question}\label{question}

En una entrevista para una documental de televisi?n, una piloto de un
helic?ptero de acrobacias comenta en detalle c?mo fue su presentaci?n:

``Durante los primeros 25 segundos aument? la altura de manera constante
y, luego, realic? una pirueta en la que descend? de manera constante
durante 10 segundos. Despu?s del descenso, comenc? a girar en torno al
eje del helic?ptero durante 25 segundos, manteniendo la misma altura.
Para finalizar, aterric? reduciendo la altura de manera constante
durante 20 segundos''.

?Cu?l de las siguientes gr?ficas representa de forma CORRECTA la
informaci?n dada por la piloto?

\subsection{Answerlist}\label{answerlist}

\begin{itemize}
\item
  \begin{figure}
  \centering
  \includegraphics{/tmp/RtmpmpGi9u/grafica_opcion_1.png}
  \caption{Opci?n A}
  \end{figure}
\item
  \begin{figure}
  \centering
  \includegraphics{/tmp/RtmpmpGi9u/grafica_opcion_2.png}
  \caption{Opci?n B}
  \end{figure}
\item
  \begin{figure}
  \centering
  \includegraphics{/tmp/RtmpmpGi9u/grafica_opcion_3.png}
  \caption{Opci?n C}
  \end{figure}
\item
  \begin{figure}
  \centering
  \includegraphics{/tmp/RtmpmpGi9u/grafica_opcion_4.png}
  \caption{Opci?n D}
  \end{figure}
\end{itemize}

\section{Solution}\label{solution}

Debemos analizar las distintas partes del vuelo seg?n la descripci?n
dada por la piloto:

\begin{enumerate}
\def\labelenumi{\arabic{enumi}.}
\tightlist
\item
  \textbf{Primera fase (0-25 segundos)}: Aumento constante de altura
  (l?nea recta ascendente).
\item
  \textbf{Segunda fase (25-35 segundos)}: Pirueta con descenso constante
  de altura (l?nea recta descendente).
\item
  \textbf{Tercera fase (35-50 segundos)}: Giro manteniendo la misma
  altura (l?nea horizontal).
\item
  \textbf{Cuarta fase (50-95 segundos)}: Aterrizaje con descenso
  constante (l?nea recta descendente).
\end{enumerate}

Analizando las opciones:

\textbf{Opci?n A:} - Subida constante hasta 25 segundos  - Descenso
constante por 10 segundos  - Mantiene altura constante durante el giro
por 25 segundos  - Descenso constante final por 20 segundos 

\textbf{Opci?n B:} - Subida constante hasta 25 segundos  - Ca?da brusca
(casi vertical)  - Deber?a ser un descenso constante - Mantiene altura
constante luego del descenso  - Descenso constante final 

\textbf{Opci?n C:} - Subida constante hasta 25 segundos  - Sube y luego
baja (formando un pico)  - Deber?a ser solo un descenso constante -
Mantiene altura constante  - Descenso constante final 

\textbf{Opci?n D:} - Subida constante hasta 25 segundos  - Baja y luego
vuelve a subir  - Deber?a ser solo un descenso constante - Mantiene
altura constante  - Descenso constante final 

Por lo tanto, la gr?fica que representa correctamente toda la
descripci?n del vuelo es la \textbf{Opci?n A}.

\subsection{Answerlist}\label{answerlist-1}

\begin{itemize}
\tightlist
\item
  Este es incorrecto. Esta gr?fica no representa correctamente las fases
  del vuelo, especialmente en la parte de la pirueta.
\item
  Este es correcto. Esta gr?fica muestra correctamente todas las fases
  del vuelo con los tiempos adecuados.
\item
  Este es incorrecto. Esta gr?fica no representa correctamente la fase
  de la pirueta, pues muestra un ascenso y luego un descenso, cuando
  deber?a ser solo un descenso constante.
\item
  Este es incorrecto. Esta gr?fica no representa correctamente la fase
  de la pirueta, pues muestra un descenso seguido de un ascenso, cuando
  deber?a ser solo un descenso constante.
\end{itemize}

\section{Meta-information}\label{meta-information}

exname: Geometria\_Vuelo\_Acrobacias extype: schoice exsolution: 0100
exshuffle: TRUE

\end{document}
